\section{Stability Analysis of SGD}
\label{sec:background}

Intuitively, smooth optimization occurs when updates reinforce each other
and gradually shift parameters toward a minimum of the loss. Larger steps
can overshoot, oscillate, and partially undo preceding parameter changes.
Understanding when such behavior remains productive, and when it produces
instability, requires examining how successive updates interact. Two factors
are central to this analysis: the \textit{local curvature} of the loss
landscape and the \textit{stochasticity} introduced by mini-batch sampling.

\subsection{Local Curvature of the Loss Function and Stability}

Training is a discrete-time dynamical system whose local response depends
on curvature. A step that produces smooth descent in a shallow direction
can overshoot in a sharper one. The Hessian quantifies this directional variation. Classical
analyses~\cite{polyak1964methods, bertsekas1997nonlinear} connect its
eigenvalues to the learning rate and local stability.

Let $\theta_t\in\mathbb{R}^p$ denote the parameters after update $t$ and
$g_t$ a mini-batch gradient evaluated at $\theta_{t-1}$. For a constant
learning rate $\eta>0$, plain SGD follows
\begin{equation}
  \theta_t=\theta_{t-1}-\eta g_t.
  \label{eq:sgd_map_stability}
\end{equation}
For full-batch GD, $g_t=\nabla\mathcal L(\theta_{t-1})$. With
$H_t:=\nabla^2\mathcal L(\theta_{t-1})$, the Jacobian of its update map is
\begin{equation}
  J_t=I-\eta H_t.
  \label{eq:jacobian_sgd}
\end{equation}
On a fixed quadratic objective, a perturbation along a Hessian eigendirection
is multiplied by $1-\eta\lambda_i$ at each step. For a positive-definite
Hessian, all components contract precisely when
\begin{equation}
  \max_i|1-\eta\lambda_i|<1
  \quad\Longleftrightarrow\quad
  \eta\lambda_{\max}(H_t)<2.
  \label{eq:quad_stability_condition}
\end{equation}
For $0<\eta\lambda_i<1$, contraction preserves the component's sign.
For $1<\eta\lambda_i<2$, its sign alternates while its magnitude shrinks.
Above two, the component alternates and grows. These cases distinguish
smooth approach, convergent oscillation, and instability of the fixed
quadratic dynamics.

Neural-network Hessians are strongly anisotropic: a few outlier directions
can be much sharper than the spectral bulk~\cite{sagun2017hessian,
ghorbani2019density}. Monitoring $\lambda_{\max}(H_t)$ therefore provides
valuable stability information, but doing so throughout training is
expensive. Even matrix-free eigensolvers require repeated Hessian--vector
products beyond the ordinary gradient computation. An online diagnostic
that recovers curvature information from those existing gradients would
make this geometric analysis accessible during routine training.

\subsection{Mini-Batch Stochasticity}

Mini-batch SGD adds a sampling-dependent component to the full-gradient
direction~\cite{mandt2017sgd, hu2019diffusion, jastrzebski2018three}.
For independent uniform batches of $b$ examples sampled without replacement
from $N>1$ examples,
\begin{equation}
\begin{aligned}
g_t&=\nabla\mathcal L(\theta_{t-1})+\xi_t,\\
\mathbb E[\xi_t\mid\theta_{t-1}]&=0,\\
\operatorname{Cov}(\xi_t\mid\theta_{t-1})
  &=\frac{N-b}{b(N-1)}\Sigma(\theta_{t-1}).
\end{aligned}
\label{eq:sgd_noise_decomposition}
\end{equation}
Here $\Sigma$ is the covariance of per-example gradients. Writing
$h_i:=\nabla\ell_i(\theta_{t-1})$ and
$\mu:=N^{-1}\sum_i h_i=\nabla\mathcal L(\theta_{t-1})$,
\begin{equation}
  \Sigma(\theta_{t-1})
  :=\frac{1}{N}\sum_{i=1}^{N}(h_i-\mu)(h_i-\mu)^\top.
  \label{eq:grad_cov_def}
\end{equation}
Its trace measures gradient variability, while its directional structure
describes how sampling perturbs the trajectory. Exact evaluation requires
per-example information absent from the aggregated training gradient.
For an inexpensive online diagnostic, the practical question is how this
sampling variability changes the alignment of the gradients already
being computed.

\subsection{Curvature, Noise and the Edge of Stability}

Learning rate and batch size also shape the geometry encountered during
training. The break-even framework of Jastrzebski et al.~\cite{jastrzebski2020break}
shows how their early values influence later curvature and gradient
variability. This motivates examining curvature and sampling together
along the trajectory.

The Edge-of-Stability (EoS) phenomenon makes the distinction between local
instability and sustained convergence especially clear. Full-batch GD
often sharpens progressively until $\lambda_{\max}(H_t)$ fluctuates near
or above $2/\eta$, while loss continues to decrease over longer
timescales~\cite{cohen2021edge, arora2022understanding}. A fixed quadratic
with the same sharpness would be unstable, but a neural network changes
its curvature as it moves. Self-stabilization analyses explain how growth
along a sharp direction can induce a reduction in curvature, limiting the
oscillation~\cite{damian2023self}. More recently, central flows describe
the time-averaged trajectory together with the covariance of its rapid
oscillations~\cite{cohen2025central}.

These results explain how short-term reversals can accompany longer-term
progress. They also make the interaction between successive updates
central to understanding convergence: the trajectory can repeatedly
reverse in sharp directions while continuing to descend elsewhere.

\subsection{Practical Stability Analysis in Large-Model Pretraining}
\label{subsec:practical-stability}

Practitioners address spikes through monitoring, recovery, and preventive
changes to the training recipe. PaLM recovered from spikes by returning
to an earlier checkpoint and skipping batches~\cite{chowdhery2022palm}.
OPT combined checkpoint recovery with learning-rate reduction and a
tighter gradient-clipping threshold~\cite{zhang2022opt}. These interventions
act on both the magnitude of the update and the state from which training
continues.

Recent open-weight reports also emphasize prevention. OLMo~2 tracks
spikes in loss and gradient norm while evaluating initialization,
normalization, and optimizer choices~\cite{olmo2025furious}. Kimi~K2's
MuonClip optimizer uses attention-logit-based rescaling of query and key
weights and reports training without loss spikes~\cite{kimi2025k2}.

Gradient norms are thus already useful monitoring signals. Their
relationship to stability becomes clearer when magnitude is combined
with direction: a growing gradient can reinforce or oppose its
predecessor. We develop this connection through constructive and
destructive regimes, showing how consecutive gradients estimate
directional curvature and reveal the conditions for a local loss
increase. This provides a geometric interpretation of signals available
during training, using one retained gradient and no Hessian computation.
